%% =====================================================================
%%  Corrigés — Unit 8 « Home sweet home »
%% =====================================================================
\corrigesUnite{8}{Home sweet home}

\begin{corrige}{1}
\begin{reponses}[3]
\item \rep{bedroom}
\item \rep{kitchen}
\item \rep{bathroom}
\item \rep{living room}
\item \rep{garden}
\end{reponses}
\end{corrige}

\begin{corrige}{2}
\begin{reponses}[3]
\item \rep{there's}
\item \rep{There are}
\item \rep{There isn't}
\item \rep{There's}
\item \rep{There are}
\item \rep{There aren't}
\end{reponses}
\pourquoi{\en{there's} (= \en{there is}) devant un nom singulier, \en{there are} devant un pluriel.
Au numéro 6, la négation au pluriel s'accompagne de \en{any}.}
\end{corrige}

\begin{corrige}{3}
\begin{reponses}[2]
\item \rep{Yes, there is.}
\item \rep{No, there isn't.}
\item \rep{Yes, there are.}
\item \rep{Yes, there is.}
\item \rep{Yes, there are.}
\item \rep{No, there isn't.}
\end{reponses}
\pourquoi{la réponse courte reprend \en{there} + \en{be}. Au numéro 6, Léa parle d'un petit
salon et d'une cuisine, mais l'appartement est petit : la cuisine n'est pas grande.}
\piege{\en{Yes, there's.} : comme toujours, pas de contraction dans une réponse courte affirmative.}
\end{corrige}

\begin{corrige}{4}
\begin{reponses}[3]
\item \rep{on}
\item \rep{under}
\item \rep{in}
\item \rep{behind}
\item \rep{between}
\item \rep{next to}
\end{reponses}
\retiens{\en{between} = entre deux choses ; \en{behind} = derrière (on voit à peine la balle) ;
\en{in front of} = devant.}
\end{corrige}

\begin{corrige}{5}
\begin{reponses}[3][french]
\item \rep{F}
\item \rep{F}
\item \rep{V}
\item \rep{F}
\item \rep{V}
\item \rep{V}
\end{reponses}
\pourquoi{c'est un \en{one-bedroom apartment} : une seule chambre. L'annonce précise \en{There isn't
a washing machine} et \en{There isn't an elevator}.}
\end{corrige}

\begin{corrige}{6}
\begin{reponses}[1]
\item \rep{There are three bedrooms.}
\item \rep{There's a garden behind the house.}
\item \rep{My keys are next to the lamp.}
\item \rep{Are there any plants?}
\item \rep{The phone is in my bag.}
\end{reponses}
\pourquoi{phrase 2 : « il y a » ne se traduit jamais par \en{it has}. Phrase 5 : à l'intérieur d'un
sac, on dit \en{in}.}
\end{corrige}

\begin{corrige}{7}
\begin{reponses}[4]
\item \rep{three}
\item \rep{sink}
\item \rep{thank}
\item \rep{sin}
\item \rep{math}
\item \rep{pass}
\item \rep{thought}
\item \rep{boat}
\end{reponses}
\pourquoi{le son \phon{θ} est un léger sifflement « plat », produit entre les dents ; \phon{s} est plus
aigu, \phon{t} plus sec. Texte de la piste page~\pageref{trans-8.10}.}
\end{corrige}

\begin{corrige}{8}
\begin{reponses}[2]
\item \rep{on his desk, next to the computer}
\item \rep{behind the door}
\item \rep{under his papers}
\item \rep{in his hand}
\end{reponses}
\pourquoi{Emma dit \en{your desk}, \en{your papers}, \en{your hand} en parlant à Jack ; en le
décrivant, on dit \en{his desk}, \en{his papers}, \en{his hand}.}
\end{corrige}

\begin{corrige}{9}
\begin{reponses}[4]
\item \rep{are}
\item \rep{There's}
\item \rep{are}
\item \rep{on}
\item \rep{next}
\item \rep{behind}
\item \rep{is}
\item \rep{opposite}
\end{reponses}
\retiens{\en{There's no elevator} = \en{There isn't an elevator} : les deux tournures sont correctes.}
\end{corrige}

\begin{corrige}{10}
\begin{reponses}[1]
\item \rep{There's a kitchen and there are two bedrooms.}
\item \rep{There isn't an elevator.} \textit{ou} \rep{There's no elevator.}
\item \rep{Is there a bathroom? — Yes, there is.}
\item \rep{Where are my keys? — They're on the table.}
\item \rep{The cat is under the bed.}
\item \rep{The laundromat is on the corner.}
\item \rep{I live in an apartment on the third floor.}
\item \rep{My phone is in my bag, next to my glasses.}
\end{reponses}
\pourquoi{phrase 7 : le « deuxième étage » français est le \en{third floor} américain, puisque le
rez-de-chaussée est le \en{first floor}. Au Royaume-Uni, on dirait \en{the second floor}.}
\end{corrige}

\begin{corrige}{11}
Réponse modèle :
\begin{anglais}
\begin{tcolorbox}[otBase,colback=otGreenTint,colframe=otGreenTint,fontupper=\fontsize{10.4}{14}\selectfont]
I live in an apartment in Lyon. It's on the fourth floor. There's a living room, a kitchen, two
bedrooms and a bathroom. There's a balcony, but there isn't a garden. There's an elevator. My
favorite room is the kitchen: there's a big table next to the window.
\end{tcolorbox}
\end{anglais}
\retiens{\en{my favorite room} = « ma pièce préférée » (britannique : \en{favourite}).}
\end{corrige}

\begin{corrige}{12}
Réponses modèles :
\begin{reponses}[1]
\item \rep{I live in an apartment.} \textit{ou} \rep{I live in a house.}
\item \rep{Yes, there is.} \textit{ou} \rep{No, there isn't.}
\item \rep{There are two bedrooms.} \textit{ou} \rep{There's one bedroom.}
\item \rep{It's next to the living room.}
\item \rep{Thank you!}
\end{reponses}
\end{corrige}

\begin{corrigetest}
\begin{reponses}[5][french]
\item \rep{b}
\item \rep{a}
\item \rep{c}
\item \rep{b}
\item \rep{a}
\item \rep{c}
\item \rep{b}
\item \rep{a}
\item \rep{c}
\item \rep{b}
\end{reponses}
\pourquoi{question 1 : \en{there are} devant un pluriel. Question 10 : \en{glasses} est un pluriel :
\en{Where are…?}}
\end{corrigetest}
